% Generado por scripts/cap10_gen.py desde documentation/legacy. No editar a mano.
{\small\setlength{\tabcolsep}{4pt}\renewcommand{\arraystretch}{1.15}
\begin{longtable}{|>{\raggedright\arraybackslash}p{4.2cm}|>{\raggedright\arraybackslash}p{11.1cm}|}
\hline
\textbf{Término} & \textbf{Definición} \\ \hline\endhead
\textbf{Campaña (Campaign)} & Partida de rol de largo plazo gestionada por un DM, con jugadores, personajes y sesiones asociadas. \\ \hline
\textbf{DM (Dungeon Master)} & Usuario que dirige la campaña; rol con permisos elevados dentro de esa campaña. \\ \hline
\textbf{Player} & Usuario miembro de una campaña que participa con uno o más personajes. \\ \hline
\textbf{Sesión (Session)} & Instancia de juego dentro de una campaña, con estado propio y notas. \\ \hline
\textbf{Personaje (Character)} & Ficha de personaje jugable perteneciente a un jugador dentro de una campaña. \\ \hline
\textbf{Enemigo (Enemy)} & Criatura/NPC antagonista gestionado por el DM dentro de una campaña. \\ \hline
\textbf{Evento (Event / SessionEvent)} & Registro de un suceso mecánico dentro de una sesión (daño, curación, XP, loot, custom) que afecta a personajes. \\ \hline
\textbf{RuleSystem} & Sistema de reglas base de un juego de rol de mesa (p. ej. "D\&D 5e"), independiente de cualquier campaña. \\ \hline
\textbf{CampaignRuleSet} & Instancia de un \texttt{RuleSystem} asociada a una campaña puntual, capaz de contener personalizaciones del DM. \\ \hline
\textbf{RuleOverride} & Personalización concreta (ban o modificación de estadística) aplicada sobre un \texttt{CampaignRuleSet}. \\ \hline
\textbf{Bitácora de Sesión (SessionLog)} & Entrada de historial narrativo y cronológico de una sesión (\texttt{combat}, \texttt{general} o \texttt{action}); las entradas \texttt{action} se generan automáticamente al registrar un \texttt{SessionEvent}. \\ \hline
\textbf{Diálogo de NPC} & Texto de diálogo generado para un NPC no planificado, mediante un servicio externo de IA. \\ \hline
\textbf{Propuesta de fecha (DateProposal)} & Fecha y horario candidatos que el DM propone para una sesión. \\ \hline
\textbf{Voto de fecha (DateVote)} & Respuesta de un jugador (sí / no / tal vez) a una propuesta de fecha. \\ \hline
\textbf{Notas de sesión} & Dos campos de texto plano por sesión: privadas (solo DM) y públicas (visibles para jugadores). \\ \hline
\end{longtable}}
